\section{Valuations on a primitive conductor shell}
\label{sec:local-shell}

We prove Theorem~\ref{thm:main} by treating a finite Laurent polynomial before
specializing to graph Laplacians. Write
\[
 F(X)=p^g\sum_{i=1}^m c_i X^{a_i},\qquad
 a_i\in\Z^d,\quad c_i\in\Z\setminus\{0\},\quad
 \min_i v_p(c_i)=0 .
\]
Equal exponents have been collected. Fix $h\ge1$, $Q=p^{h-1}$,
$e=(p-1)Q$, and a primitive $p^h$-th root $\zeta$.
A point of joint conductor $p^h$ is $\zeta^u=(\zeta^{u_1},\ldots,
\zeta^{u_d})$, where at least one entry of $u\bmod p^h$ is a unit.
The simultaneous unit action on such tuples is free. Its orbits
have size $e$, and the valuation is constant on an orbit since
$\Q_p(\zeta)$ has a unique normalized $p$-adic valuation.
Thus a single representative contributes $e v_p(F(\zeta^u))$ to the
sum over the whole shell.

For each monomial choose $0\le n_i<pQ$ congruent to $a_i\cdot u$
modulo $pQ$, and write $n_i=r_i+Qb_i$ with $0\le r_i<Q$.
Set $A_{rb}=\sum_{i:(r_i,b_i)=(r,b)}c_i$. The unique representative
of the normalized evaluation $p^{-g}F(\zeta^u)$ in the power basis
$1,\zeta,\ldots,\zeta^{e-1}$ is $R(\zeta)$, where
\[
 R(Y)=\sum_{r=0}^{Q-1}\sum_{b=0}^{p-2}
 C_{rb}Y^{r+Qb},\qquad C_{rb}=A_{rb}-A_{r,p-1}.
\]
Indeed $\Phi_{p^h}(Y)=1+Y^Q+\cdots+Y^{(p-1)Q}$ eliminates the
last block of exponents. In particular $R=0$ detects exact vanishing.

\begin{lemma}[Integral orbit valuation]
\label{lem:orbit-val}
If $R\ne0$, let $\mu=\min_{C_{rb}\ne0} v_p(C_{rb})$.
Let $j$ be the first index whose coefficient in
$p^{-\mu}R(1+T)$ is nonzero modulo $p$. Then $0\le j<e$ and
\[
 e v_p(F(\zeta^u))=e(g+\mu)+j .
\]
\end{lemma}

\begin{proof}
Changing variables from $Y$ to $1+T$ preserves nonzero polynomials
over $\F_p$, so $j$ exists. The element $\pi=\zeta-1$ has
$v_p(\pi)=1/e$. In $p^{-\mu}R(1+\pi)$ the term of degree $j$
has valuation $j/e<1$. Every earlier nonzero coefficient is divisible
by $p$ and contributes valuation at least $1$, and each later term
has valuation strictly greater than $j/e$. The unique term of least
valuation gives the assertion, also when $p=2$ and $h=1$.
\end{proof}

The integer content is uniformly bounded: the sums defining the two terms in $C_{rb}$ use disjoint subsets of the coefficients, so $p^\mu\le\sum_i|c_i|$ when $R\ne0$. For a small example take $p=2,h=2$ and the square-grid symbol at $u=(1,1)$. The exponent polynomial $4-2X-2X^3$ reduces to $4$ modulo $X^2+1$. Thus $g=0,\mu=2,j=0$ and the orbit contributes $4$. Computing the content before collecting the terms would give the wrong answer.

\section{Computing all conductors with finitely many states}
\label{sec:lucas-machine}

Normalize each representative by making its first unit coordinate
equal to $1$ modulo $p^h$. In the chart with pivot $i_0$, lower-indexed
coordinates have units digit zero. The pivot's units digit is one, and
its other digits are zero. Every other digit is free. Reading the
coordinate vectors from the least significant end therefore
enumerates each orbit exactly once. At every stage there are
two operations on a new digit. Treating it as the top digit
emits a valuation at that conductor. Treating it as an ordinary digit
prepares the next stage.

After $k$ ordinary digits let $Q=p^k$, let $w\in\{0,\ldots,Q-1\}^d$
be the prefix, and decompose
\[
 a_i\cdot w=r_i+Q s_i,\qquad 0\le r_i<Q .
\]
Only the carries $s_i$ and a small amount of information about the $r_i$
are needed. For $0\le j<Q$ define a row vector over $\F_p$ by
\[
 B_j=\left(\binom{r_1}{j},\ldots,\binom{r_m}{j}\right)\pmod p .
\]
Process these rows in order of increasing $j$, retaining a row precisely
when it is outside the span of the earlier retained rows. Store the
ordered list of retained vectors and keep their integer indices in
separate registers. There are at most $m$ such registers.

\begin{lemma}[What the retained rows determine]
\label{lem:lucas-basis}
Two columns agree on the retained rows if and only if their residues
$r_i$ agree. For any linear functional on the rows, its first nonzero
value, if one exists, occurs on a retained row. The retained list after
one ordinary digit is computable from the previous list and the carries.
\end{lemma}

\begin{proof}
The full column associated to $r$ is the coefficient vector of
$(1+T)^r\in\F_p[T]$. Distinct $r<Q$ give different degrees. Agreement
on a basis implies agreement on all rows. A discarded row lies
in the span of earlier rows, so it cannot be the first detected by a
linear functional.

For a new coordinate digit vector $v\in\{0,\ldots,p-1\}^d$, perform
Euclidean division $s_i+a_i\cdot v=b_i+p s'_i$, $0\le b_i<p$.
The new residue is $r_i+Qb_i$ and the new carry is $s'_i$.
Over $\F_p$ one has $(1+T)^Q=1+T^Q$, hence
\[
 \binom{r_i+Qb_i}{j+Qt}
 \equiv \binom{r_i}{j}\binom{b_i}{t}\pmod p
 \qquad(0\le j<Q,\;0\le t<p).
\]
For each $t=0,\ldots,p-1$ in order, multiply every old retained row
coordinatewise by $(\binom{b_i}{t})_i$. Feed these candidates, in old
row order within the block, to the same greedy independence test.
The omitted old rows are combinations of earlier rows in their own
block after this multiplication. Inserting them into the scan would
therefore leave all decisions unchanged. A retained candidate coming
from an old index $j$ has new index $j+Qt$.
\end{proof}

\paragraph{Reading the top digit.}
Use the carries to compute $b_i$ as above. Partition the columns
according to their entries on the stored basis.
Lemma~\ref{lem:lucas-basis} identifies these groups with the occupied
values of $r$. Compute the integers $C_{rb}$ in each group, test for
vanishing, and then compute $\mu$ by Lemma~\ref{lem:orbit-val}.
Division by $p^\mu$ takes place \emph{after} the integer sums and
subtractions within each group. In top block $t=0,\ldots,p-2$ put
\[
 w_r(t)=\sum_{b=0}^{p-2}(C_{rb}/p^\mu)\binom bt\pmod p .
\]
At index $j_0+Qt$, the coefficient to be tested modulo $p$ is
$\sum_r w_r(t)\binom{r}{j_0}$. This is a linear functional on
the old rows, using one coordinate per group. Testing the stored rows
in order finds its earliest nonzero value. Stop at the first successful
block. If the successful old index is $j_0$, the orbit contribution is
\[
 \big((p-1)(g+\mu)+t\big)Q+j_0 .
\]
All choices except the numerical values of $Q,j_0$ depend on the finite
state and this last digit.

\paragraph{A reduction of the state space.}
After processing $k$ ordinary digits, test the normalized coefficient
functional $B_j\mapsto\sum_i c_i\binom{r_i}{j}$. Suppose its first
nonzero value is at $j_*<p^k$. For every subsequent conductor the
ordinary exponent polynomial at $1+T$, reduced modulo $(p,T^{p^k})$,
equals $\sum_i c_i(1+T)^{r_i}$ in this quotient.
Indeed further exponent digits multiply these terms by powers of
$1+T^{p^k}$. Final cyclotomic reduction preserves the same truncation,
because $\Phi_{p^h}(1+T)\equiv T^{\varphi(p^h)}\bmod p$ and
$\varphi(p^h)\ge p^k$ for $h>k$. Consequently every continuation has
$\mu=0$ and $j=j_*$, including arbitrary choices for its final digit.
The implementation stores only $Q$ and $j_*$ in a new state with the
same chart. Future ordinary digits replace $Q$ by $pQ$ and leave
$j_*$ fixed.

\paragraph{Termination and weights.}
The control state consists of the pivot, a first-digit flag, the
carries, and the ordered list of row vectors. The register values are
stored separately. If $H_i=\sum_\ell |a_{i\ell}|$, then
$|s_i|=|\lfloor a_i\cdot w/Q\rfloor|\le H_i$ for every prefix.
An explicit upper bound on the number of control states, including the
reduced states, is
\[
 N=d+2d\left(\prod_{i=1}^m(2H_i+1)\right)\sum_{s=0}^m p^{ms}.
\]
A breadth-first search enumerates the reachable states and tests every
allowed final digit. It therefore also decides whether a nontrivial
$p$-power torsion zero exists. One can reject such inputs or, for a
sum taken only over nonzero evaluations, give a zero evaluation
weight zero while continuing the search. Rejection is appropriate
for an ordinary product.

Each ordinary state has registers $(Q,j_1,\ldots,j_s)$ with $s\le m$.
Each ordinary digit induces an integer linear map between registers by
$Q'=pQ$, $j'_a=j_b+tQ$. Summing the maps for digits with the same target
gives the corresponding block of a matrix $M$. In each state the sum
of the output forms over its legal top digits gives the corresponding
block of a row $f$. Initially each chart has zero carries, the single
row $(1,\ldots,1)$, and registers $(1,0)$. Let $u$ be the concatenated
initial vector, zero outside these initial states. Matrix iteration
adds the registers of all prefixes at each state. By induction on
prefix length and the unique chart representation, the shell sum at
conductor $p^h$ is exactly
\[
 b_h=fM^{h-1}u .
\]
The matrix dimension $L$ is the sum of the register counts of the
reachable states. In particular $L\le(m+1)N$. The implementation
uses the dimension obtained from its search.

\section{Spanning trees and the first level of the eventual law}
\label{sec:exact-gf}

Let $G$ be a finite connected undirected multigraph with integral
voltages $a_e\in\Z^d$, one orientation per edge. Its Laurent Laplacian
adds $1$ to both endpoint diagonals and subtracts $X^{a_e},X^{-a_e}$
in the two corresponding off-diagonal positions, interpreting all four
entries at the same position for a loop. Take $F$ to be its determinant.
The finite lifts have deck group $(\Z/p^n\Z)^d$ and we assume they
are connected. Character decomposition of the lifted Laplacian
gives one block per character. The trivial block is the base Laplacian.
The matrix-tree theorem applied to the base and lift cancels their
common vertex-count factor and leaves the deck-group order in the
denominator. Hence, writing $v_n=v_p(\kappa_n)$ and
$v_0=v_p(\kappa(G))$,
\[
 v_n=v_0-dn+\sum_{h=1}^n b_h .
\]
The base cofactor and Laurent determinant are explicit finite input
calculations. The same argument works for graphs with leaves and for
cycles. A connected base has a connected lift at a given level exactly
when every nontrivial character block there is invertible, because
the multiplicity of the zero eigenvalue counts components.

With $z$ a formal variable, summing the last display yields
\begin{equation}
 \sum_{n\ge0}v_n z^n=
 \frac{v_0}{1-z}-\frac{dz}{(1-z)^2}
 +\frac{z}{1-z}\,f(I-zM)^{-1}u .
 \label{eq:machine-resolvent}
\end{equation}
This identity computes the entire sequence. If
$\chi_M(x)=x^L+c_1x^{L-1}+\cdots+c_L$, then
$\det(I-zM)=1+c_1z+\cdots+c_Lz^L$. The adjugate shows that the
numerator of $f(I-zM)^{-1}u$ over this determinant has degree at most
$L-1$. Multiplying that denominator by the first $L$ shell coefficients
of the series and keeping degrees below $L$ gives its exact numerator,
including when $M$ is singular. This is the characteristic-polynomial
procedure used in the accompanying implementation.

For the Greenberg polynomial put
\[
 D(z)=(1-p^d z)\prod_{i=0}^{d-1}(1-p^i z)^2 .
\]
An eventual law $v_n=P(p^n,n)$ with $\deg P\le d$ and
$\deg_Y P(X,Y)\le1$ implies that the reduced denominator of
\eqref{eq:machine-resolvent} divides $D(z)$. Conversely, whenever it
divides $D(z)$, polynomial division and partial fractions give
\[
 \sum_{n\ge0}v_nz^n
 =H(z)+\sum_{i=0}^d\left(\frac{A_i}{1-p^i z}
            +\frac{B_i}{(1-p^i z)^2}\right),\qquad B_d=0 .
\]
Thus $P(X,Y)=\sum_i[A_i+(Y+1)B_i]X^i$ and
$v_n-P(p^n,n)=[z^n]H$. If $H=0$ the least onset is zero.
Otherwise it is $1+\deg H$. Uniqueness follows since the largest
exponential base in a nonzero sum $\sum_i(a_i+b_i n)p^{in}$ dominates
its lower bases on a tail. DuBose--Vallières guarantee the denominator
shape for base graphs with no leaves and nonzero Euler characteristic.
To pass to an arbitrary connected base admitting such a connected tower,
repeatedly remove pendant edges. Their lifts are pendant edges and belong
to every spanning tree. If the remaining core is a cycle, its fundamental
group is cyclic, so a connected tower with deck groups as above has
$d=1$. A cycle of length $b$ then lifts to a cycle of length $bp^n$,
with $v_n=v_p(b)+n$ (including a single loop or a double edge).
A tree admits no connected nontrivial finite cover. These cases exhaust
the missing base graphs. There is also a numerical upper bound on the required conductor precision. Formula~\eqref{eq:machine-resolvent} has numerator and denominator degrees at most $L+2$. Consequently $\deg H\le L+2$ if $H\ne0$, and the least onset is at most $L+3$. The characteristic-polynomial computation itself requires only the first $L$ shell sums after the reachable matrix has been built. Computed divisibility
also certifies the shape directly in each example. For the general
Laurent sum, the characteristic-polynomial procedure produces its
exact rational series independently of that existence theorem.

The integer-voltage formulation also covers rational voltages in
$\Q\cap\Z_p$. Multiplying every voltage by a common denominator prime
to $p$ induces compatible deck-group automorphisms at every level,
so the finite graphs are isomorphic to those with the cleared
integral voltages.
